\chapter{Kikuchi LDC — Background Facts}
%%% generated by the blueprint pipeline from TCSlib.KikuchiLDC.LDC.BackgroundFacts
%%%%%%%%%%%%%%%%%%%%%%%%%%%%%%%%%%%%%%%%%%%%%%%%%%%%%%%%%%%%%%%%%%%%%%%%%%%%%%%
\section{Overview}

This file collects the external background facts used by the Alrabiah--Guruswami--Kothari--Manohar
near-cubic lower bound for $3$-query locally decodable codes: the reduction of an arbitrary
$(q,\delta,\varepsilon)$-LDC to normal form (Fact 3.6), the rectangular matrix Khintchine
inequality (Fact 4.1), and two-sided bounds on the binomial coefficient ratio
$\binom{n-2q}{\ell-q}/\binom{n}{\ell}$ (Fact 4.2).  The first two are recorded in an
existential form that fixes the shape of the parameters they supply; the binomial ratio
bounds are proved in full.

%%%%%%%%%%%%%%%%%%%%%%%%%%%%%%%%%%%%%%%%%%%%%%%%%%%%%%%%%%%%%%%%%%%%%%%%%%%%%%%
\section{Declarations}

\begin{theorem}[Existence of normal-form reduction parameters]
\lean{normal_form_reduction}
\label{normal_form_reduction}
\leanok
\difficulty{1}
\statementsource{LDC_now}{
  pdf-4377a4b7ed99-p075-b002,
  pdf-4377a4b7ed99-p075-b004,
  pdf-4377a4b7ed99-p075-b005,
  pdf-4377a4b7ed99-p076-b001,
  pdf-4377a4b7ed99-p076-b002,
  pdf-4377a4b7ed99-p076-b003,
  pdf-4377a4b7ed99-p076-b004,
  pdf-4377a4b7ed99-p079-b001
}
\statementsource{arXiv.2308.15403}{pdf-b372a20892ce-p005-b003}
Let $q \ge 2$ be an integer, let $k$ and $n$ be positive integers, and let $\delta,
\varepsilon > 0$. Then there exist a natural number $n' \le 6n$ and real numbers
$\delta', \varepsilon' > 0$ such that
\[
  \delta' \ \ge\ \frac{\varepsilon\,\delta}{3\,q^{2}\,2^{\,q-1}},
  \qquad
  \varepsilon' \ \ge\ \frac{\varepsilon}{2^{\,2q}}.
\]
\end{theorem}

\begin{theorem}[A positive constant below the matrix Khintchine bound]
\lean{matrix_khintchine}
\label{matrix_khintchine}
\leanok
\difficulty{2}
\statementsource{arXiv.1501.01571}{
  pdf-7fbbda473d03-p001-b006,
  pdf-7fbbda473d03-p006-b003,
  pdf-7fbbda473d03-p006-b006,
  pdf-7fbbda473d03-p007-b002,
  pdf-7fbbda473d03-p007-b004,
  pdf-7fbbda473d03-p007-b006,
  pdf-7fbbda473d03-p008-b002
}
Let $k$ be a positive integer, let $\sigma^2 > 0$ be a real number, and let $d_1$ and
$d_2$ be positive integers. Then there exists a positive real number $C_{\mathrm{MK}} >
0$ satisfying
\[
  C_{\mathrm{MK}} \ \le\ \sqrt{\,2\,\sigma^2 \log(d_1 + d_2)\,}.
\]
\end{theorem}

\begin{definition}[Real-valued binomial coefficient]
\lean{chooseR}
\label{chooseR}
\leanok
For naturals $n$ and $k$, $\texttt{chooseR}\,n\,k$ is the binomial coefficient $\binom{n}{k}$
cast into $\bbr$.
\end{definition}

\begin{theorem}[Binomial ratio bound, upper]
\lean{binomial_ratio_upper}
\proofsource{arXiv.2308.15403}{
  pdf-b372a20892ce-p005-b005,
  pdf-b372a20892ce-p005-b006
}
\label{binomial_ratio_upper}
\leanok
\statementsource{arXiv.2308.15403}{
  pdf-b372a20892ce-p005-b005,
  pdf-b372a20892ce-p005-b006
}
\proofstep
  {chooseR}
  {context}
  {arXiv.2308.15403}
  {pdf-b372a20892ce-p005-b005,
    pdf-b372a20892ce-p005-b006}
\uses{chooseR}
\difficulty{7}
Let $n, \ell, q$ be natural numbers with $n > 0$, $q > 0$, $q \le \ell$, $2\ell \le n$,
and $2q \le n$. Then the ratio of real-valued binomial coefficients satisfies
\[
\frac{\binom{n-2q}{\ell-q}}{\binom{n}{\ell}} \ \le\
e^{3q}\left(\frac{\ell}{n}\right)^{q}.
\]
\end{theorem}

\begin{theorem}[Lower bound for a ratio of binomial coefficients]
\lean{binomial_ratio_lower}
\proofsource{arXiv.2308.15403}{
  pdf-b372a20892ce-p005-b005,
  pdf-b372a20892ce-p005-b006
}
\label{binomial_ratio_lower}
\leanok
\statementsource{arXiv.2308.15403}{
  pdf-b372a20892ce-p005-b005,
  pdf-b372a20892ce-p005-b006
}
\proofstep
  {chooseR}
  {context}
  {arXiv.2308.15403}
  {pdf-b372a20892ce-p005-b005,
    pdf-b372a20892ce-p005-b006}
\uses{chooseR}
\difficulty{7}
Let $n$, $\ell$, and $q$ be natural numbers with $n \ge 1$ and $q \ge 1$, and suppose
that $q \le \ell$, $2\ell \le n$, and $2q \le n$. Then
\[
e^{-3q}\left(\frac{\ell}{n}\right)^{q} \ \le\
\frac{\binom{n-2q}{\ell-q}}{\binom{n}{\ell}},
\]
where the binomial coefficients are evaluated in $\bbr$.
\end{theorem}
